\documentclass[12pt]{article}
\usepackage{amsmath, amssymb, amsthm}
\usepackage{mathtools}
\usepackage{hyperref}
\usepackage{listings}
\usepackage{xcolor}

\lstdefinelanguage{lean}{
  keywords={import, namespace, open, variable, def, lemma, theorem, structure, inductive, deriving, section, end},
  keywordstyle=\color{blue}\bfseries,
  ndkeywords={sorry, by, exact},
  ndkeywordstyle=\color{red}\bfseries,
  comment=[l]{--},
  commentstyle=\color{gray}\ttfamily,
  stringstyle=\color{green}\ttfamily,
  morestring=[b]",
  basicstyle=\ttfamily\small,
  breaklines=true
}

\title{BOSONIZE-LEAN: Phase 2 Formalization Blueprint}
\author{Yellow Team Architect}
\date{\today}

\begin{document}
\maketitle

\section{Phase 2 Scope Constraints}
Based on the Phase 2 directives and critiques, this phase is strictly limited to the formulation of the Bosonization Dictionary's density sector. Klein factors, Mattis-Mandelstam, and Luttinger interaction are deferred to subsequent phases. This document exactly mirrors the required Lean 4 stubs in \texttt{Bosonize/Stubs/Phase2\_Stubs.lean}.

\section{2.1 Density Operators}
We define the density modes as the Fourier transform of the normal-ordered fermion density. Normal ordering subtracts the macroscopic vacuum expectation value to keep the operator finite on the budget subspace.

\begin{lstlisting}[language=lean]
section Densities

/-- The raw density operator \rho(m) = \sum c^\dagger_{q+m} c_q -/
def density_mode (m : LambdaDual L) : Module.End Complex (FockSpace (Chirality x LambdaDual L)) := 
  -- Note to Blue Team: implement using Finset.sum over q in LambdaDual
  sorry

/-- Vacuum expectation value of the density mode -/
def density_vev (m : LambdaDual L) : Complex :=
  -- Tr( |0><0| \rho(m) )
  sorry

/-- Normal ordered density mode :\rho(m): -/
def no_density_mode (m : LambdaDual L) : Module.End Complex (FockSpace (Chirality x LambdaDual L)) :=
  density_mode m - LinearMap.smulRight (LinearMap.id) (density_vev m)

/-- The zero mode of the normal ordered density is exactly the particle number shift -/
lemma no_density_zero_eq_N : 
  no_density_mode 0 = N_op := sorry

end Densities
\end{lstlisting}

\section{2.2 Heisenberg Algebra on the Budget Subspace}
The Schwinger term arises algebraically. The commutator is evaluated strictly on the finite Energy Budget Subspace $\mathcal{B}_{K, N_{max}}$ to rigorously avoid infinite limits and domain issues.

\begin{lstlisting}[language=lean]
section HeisenbergCCR

/-- 
  The Exact Schwinger Term on the Budget Subspace.
  Requires the momentum modes m, n to be small enough such that they do not 
  excite states out of the finite band L when acting on a state in the budget K.
-/
lemma comm_no_density_modes {K Nmax : Nat} {m n : LambdaDual L} 
  (h_band : (m.val : Int) + (n.val : Int) + K < L / 2) :
  \forall (\psi : BudgetSpace L K Nmax), 
    (no_density_mode m * no_density_mode n - no_density_mode n * no_density_mode m) \psi.val = 
    (if m.val + n.val = 0 then (m.val : Complex) else 0) \bullet \psi.val := sorry

end HeisenbergCCR
\end{lstlisting}

\section{2.3 Sugawara Construction}
The construction of the bosonic kinetic energy operator. On the budget subspace, this is rigorously proven equivalent to the free linearized fermionic Hamiltonian.

\begin{lstlisting}[language=lean]
section Sugawara

/-- The bosonic kinetic energy operator -/
def sugawara_op : Module.End Complex (FockSpace (Chirality x LambdaDual L)) :=
  \sum m in (Finset.filter (fun x : LambdaDual L => (x.val : Int) > 0) Finset.univ),
    no_density_mode (-m) * no_density_mode m

/-- 
  The free fermion Hamiltonian equals the Sugawara bosonic Hamiltonian 
  on the budget subspace, up to the zero-mode energy shift.
-/
lemma sugawara_equivalence {K Nmax : Nat} (h_band : K < L / 4) :
  \forall (\psi : BudgetSpace L K Nmax),
    -- H_0 \psi = (H_sug + N(N+1)/2) \psi
    sorry

end Sugawara
\end{lstlisting}

\section{2.4 Haldane Completeness}
Proving that the set of states generated by acting with arbitrary polynomials of $\rho(-m)$ forms a complete basis for the budget subspace via integer partition counts.

\begin{lstlisting}[language=lean]
section Completeness

/-- 
  Dimension count showing that the number of linearly independent states 
  generated by \rho(-m) matches the dimension of the fixed-N budget subspace, 
  proving that the bosons completely span the fermionic degrees of freedom.
-/
lemma haldane_completeness {K N : Nat} (h_band : K < L / 4) :
  Module.finrank Complex (sorry : Submodule Complex (FockSpace (Chirality x LambdaDual L))) = 
  Nat.Partition.partitions K := sorry

end Completeness
\end{lstlisting}

\end{document}
